\documentclass[12pt]{article}
\usepackage[utf8]{inputenc}
\usepackage[T1]{fontenc}
\usepackage{hyperref}
\title{Tobelo Language Resources for K-12 Curriculum Development}
\author{Compiled by Claude Sonnet 5 (claude-sonnet-5) for Andr\'as Kornai}
\date{}
\begin{document}
\maketitle

\section{Introduction}
Tobelo (ISO 639-3 tlb) is a North Halmahera language of North Maluku, Indonesia, with roughly 30,000 speakers. Searches found general reference pages and one published sketch grammar, but no school materials. Little exists in the results, and this report says so plainly.

\section{Primary Educational Resources}
No primers, textbooks or school programmes for Tobelo turned up in the search results.

\section{Linguistic and Reference Materials}
\begin{itemize}
\item Overview of the language, its dialects and relatives: \url{https://en.wikipedia.org/wiki/Tobelo_language}
\item Omniglot entry on Tobelo writing: \url{https://www.omniglot.com/writing/tobelo.htm}
\item The main grammatical description is Gary Holton's \emph{Tobelo} (LINCOM, 2003), a sketch grammar with texts. The search results described it but gave no direct link to the book.
\end{itemize}

\section{Religious and Community Resources}
No Tobelo Bible translation or community organisation was found in the search results.

\section{Digital and Audio Resources}
\begin{itemize}
\item PARADISEC field-notes page for North Halmahera: \url{https://www.paradisec.org.au/fieldnotes/INHALM.htm}
\end{itemize}

\section{Curriculum Development Context}
There is very little ready-made material. A curriculum would have to be built from the Holton grammar and texts, and from new community-produced material.
\end{document}